\begin{center}
    \Large{\textbf{Аннотация}}
\end{center}

\begin{singlespace}
В данной работе исследуется проблема структурного прореживания нейронных сетей. Структурное прореживание представляет собой процесс удаления групп параметров из нейронной сети, например, фильтров в свёрточных нейронных сетях. Оптимальная стратегия прореживания позволяет повысить как обобщающую способность модели, так и эффективность выполнения инференса. Основная сложность структурного прореживания заключается в том, что удаление одного слоя требует исключения или согласованного изменения всех зависимых слоев.

Предложенный метод основан на анализе вычислительного графа нейронной сети и оценке распространения информации по графу. Данный подход позволяет оценивать значимость операций в вычислительном графе в условиях few-shot обучения, то есть с использованием нескольких прямых проходов модели на подмножестве данных. Для демонстрации эффективности предложенного метода проведен ряд экспериментов на датасете MNIST.
\end{singlespace}